% p.18 — الباب التاسع في الملحقات بالرباعي
% Layout (as in the lithograph): under the chapter title, one `tables` row holds the main
% table on the RIGHT (three blocks: ملحق بدحرج 1–7, ملحق بتدحرج 1–7, ملحق باحرنجم 1–3) and,
% on the LEFT, the prose column headed «قانون» (\sidetext, the full height of that table).
% Below both runs the full-width table of اِفْتَعَل → فَعَّل with the letters ص ض ط ظ د ذ ز ث و ي
% in two narrow right-hand columns, whose last two letter-rows share one «قانون» cell and one
% prose cell spanning the remaining columns.
\mspage{18}
\tablereserve{13}
\bab{الباب التاسع في الملحقات بالرباعي}
\begin{tables}
\begin{sarf}{|c|c|c|c|c|c|}
\thead{} & \thead{ثلاثى مجرّد} & \thead{ثلاثى مزيد} & \thead{ثلاثى مجرد} & \thead{ثلاثى مزيد} & \thead{ملحق بالرباعى}\\\hline
1 & فَعَلَ & فَعْلَلَ & جَلَبَ & جَلْبَبَ & دَحْرَجَ\\
2 & فَعَلَ & فَعْنَلَ & قَلَسَ & قَلْنَسَ & دَحْرَجَ\\
3 & فَعَلَ & فَوْعَلَ & جَرَبَ & جَوْرَبَ & دَحْرَجَ\\
4 & فَعَلَ & فَعْوَلَ & سَرَلَ & سَرْوَلَ & دَحْرَجَ\\
5 & فَعَلَ & فَيْعَلَ & \unclear{خَتَلَ} & \unclear{خَيْتَلَ} & دَحْرَجَ\\
6 & فَعَلَ & فَعْيَلَ & شَرَفَ & شَرْيَفَ & دَحْرَجَ\\
7 & فَعَلَ & فَعْلَى & قَلَسَ & قَلْسَى & دَحْرَجَ\\\hline
1 & فَعَلَ & تَفَعْلَلَ & جَلَبَ & تَجَلْبَبَ & تَدَحْرَجَ\\
2 & فَعَلَ & تَفَعْنَلَ & قَلَسَ & تَقَلْنَسَ & تَدَحْرَجَ\\
3 & فَعَلَ & تَفَوْعَلَ & جَرَبَ & تَجَوْرَبَ & تَدَحْرَجَ\\
4 & فَعَلَ & تَفَعْوَلَ & سَرَلَ & تَسَرْوَلَ & تَدَحْرَجَ\\
5 & فَعَلَ & تَفَيْعَلَ & \unclear{خَتَلَ} & \unclear{تَخَيْتَلَ} & تَدَحْرَجَ\\
6 & فَعَلَ & تَفَعْيَلَ & شَرَفَ & تَشَرْيَفَ & تَدَحْرَجَ\\
7 & فَعَلَ & تَفَعْلَى & قَلَسَ & تَقَلْسَى & تَدَحْرَجَ\\\hline
1 & فَعَلَ & اِفْعَنْلَلَ & قَعَسَ & اِقْعَنْسَسَ & اِحْرَنْجَمَ\\
2 & فَعَلَ & اِفْوَنْعَلَ & حَصَلَ & اِحْوَنْصَلَ & اِحْرَنْجَمَ\\
3 & فَعَلَ & اِفْعَنْلَى & سَلَقَ & اِسْلَنْقَى & اِحْرَنْجَمَ\\
\end{sarf}
% Left-hand prose column headed «قانون»:
\sidetext{0.36\linewidth}{%
{\centering\bfseries قانون\par}
الثلاثى المجرد = ما كان ماضيه على ثلاثة احرف و لم يكن فيه زائدٌ وهى ستة ابواب (1) فَتَحَ (2) ضَرَبَ (3) نَصَرَ (4) سَمِعَ (5) حَسِبَ (6) كَرُمَ
\par
الثلاثى المزيد = ما كان ماضيه على ثلاثة احرف وكان فيه حرف زائد او حرفان او ثلاثة احرف وهى اَرْبَعَةَ عَشَرَ (1) أَفْعَلَ (2) فَعَّلَ (3) فَاعَلَ (4) تَفَعَّلَ (5) تَفَاعَلَ (6) اِفْتَعَلَ (7) اِنْفَعَلَ (8) اِفْعَلَّ (9) اِسْتَفْعَلَ (10) اِفَّعَّلَ (11) اِفَّاعَلَ (12) اِفْعَالَّ (13) اِفْعَوَّلَ (14) اِفْعَوْعَلَ
\par
الرباعى المجرد = ما كان ماضيه على اربعة احرف و لم يكن فيه زائد وهو واحد (1) فَعْلَلَ مثل دَحْرَجَ
\par
الرباعى المزيد = ما كان ماضيه على اربعة احرف وكان فيه حرف زائد او حرفان و هى ثلاثة (1) تَفَعْلَلَ تَدَحْرَجَ (2) اِفْعَنْلَلَ اِحْرَنْجَمَ (3) اِفْعَلَلَّ اِقْشَعَرَّ
\par
الملحق بالرباعى = ما كان ماضيه على ثلاثة احرف وكان فيه حرف زائد او حرفان او ثلاثة احرف وكان تصريفه مثل الرباعى، فالملحق بدحرج سبعة وبتدحرج سبعة و باحرنجم ثلاثة ومعنى الالحاق الاتحاد فى المصدر
}
\end{tables}
% Bottom table. The two right-hand columns list the letters in pairs (ص/ض, ط/ظ, د/ذ, ز/ث, و/ي);
% the last two letter-rows (ز/ث, و/ي) share one «قانون» cell and one prose cell spanning the remaining columns.
\begin{sarf}{|c|c|c|c|c|c|c|c|}
ص & ض & اِفْتَعَلَ & يَفْتَعِلُ & اِفْتِعَالًا & مُفْتَعِلٌ & اِفْتَعِلْ & لَا تَفْتَعِلْ\\
ط & ظ & فَعَّلَ & يَفَعِّلُ & فِعَّالًا & مُفَعِّلٌ & فَعِّلْ & لَا تَفَعِّلْ\\
د & ذ & خَصَّمَ & يَخَصِّمُ & خِصَّامًا & مُخَصِّمٌ & خَصِّمْ & لَا تَخَصِّمْ\\\hline
\twoline{ز}{و} & \twoline{ث}{ي} & \thead{قانون} & \multicolumn{5}{p{0.55\linewidth}|}{اِنْ كَانَ حَرْفٌ مِنْ هٰذِهِ الْحُرُوْفِ بعد تاءِ اِفتعل فَيَجُوْزُ اَنْ يُبَدَّلَ اِفْتَعَلَ بِفَعَّلَ نَحْوُ فِي اِخْتَصَمَ خَصَّمَ وَتَصْرِيْفُهُ كَمَا ذُكِرَ}\\
\end{sarf}
